\subsection{可测函数的本质像}

在本节中，我们用 X 表示局部紧拓扑空间，用 \(\mu\) 表示 X 上的正测度。

定义 1. --- 设 Y 是拓扑空间。对每个从 X 到 Y 中的可测函数 g，\(\mu\)-本质像（\(g\) 的）是满足如下条件的 \(y \in Y\) 所成之子集：对每个开邻域 \(U\) （\(y\) 的），集 \(g^{-1}(\mathrm{U})\) 在 X 中不是局部 \(\mu\)-零测的（INT, IV, p. 172, § 5, no. 2，定义 3）。

引理 1. --- 设 \(g\) 是从 X 到拓扑空间 Y 中的 \(\mu\)-可测函数，并设 \(S\) 是它的 \(\mu\)-本质像。元素 \(x \in X\) 中使得 \(g(x)\) 不属于 S 者所成之集在 X 中局部 \(\mu\)-零测。

设 \(Z = g^{-1}(Y - S)\) 是所考虑的集。

首先假设 X 紧且 g 连续。在此情形，前推测度 \(g(\mu)\) 有定义（INT, V, p. 69, § 6, no. 1，定义 1），因为 \(\mu\) 是有界测度。由定义可得，\(\mu\)-本质像（\(g\) 的）是测度 \(g(\mu)\) 的支撑（参见 INT, V, p. 70, § 6, no. 2，推论 2），于是 \(\mu(Z) = g(\mu)(Y - S) = 0\) （INT, IV, p. 118, § 2, no. 2，命题 5）。

现在考虑一般情形。设 C 是 X 的紧子集，\(\varepsilon > 0\) 是实数。由于 \(g\) 可测，存在紧子集 \(C_1 \subset C\) ，使得 \(\mu(C - C_1) \leqslant \varepsilon\) 且 \(g|C_1\) 连续（INT, IV, p. 169, § 5, no. 1，命题 1）。于是我们有

\[
\mu (\mathrm{Z} \cap \mathrm{C}) \leqslant \mu (\mathrm{C} - \mathrm{C} _ {1}) + \mu (\mathrm{Z} \cap \mathrm{C} _ {1}) \leqslant \varepsilon + \mu (\mathrm{Z} \cap \mathrm{C} _ {1}).
\]

设 \(\mu_{1}\) 是 \(\mu\) 在 \(C_{1}\) 上诱导的测度（INT, IV, p. 186, § 5, no. 7，定义 4）。设 \(S_{1}\) 是 \(\mu_{1}\)-本质像（\(g|C_{1}\) 的）。我们有 \(Z\cap C_{1}\subset(g|C_{1})^{-1}(Y-S_{1})\) 。由第一种情形，集 \(Z\cap C_{1}\) 因此是 \(\mu_{1}\)-零测的，从而是 \(\mu\)-零测的（INT, IV, p. 187, § 5, no. 7，引理 2, (i)）。上面的不等式变成 \(\mu(Z\cap C)\leqslant\varepsilon\) ；由于 \(\varepsilon\) 与 C 是任意的，集 Z 局部 \(\mu\)-零测（INT, IV, p. 172, § 5, no. 2，命题 5）。

引理 2. --- 设 \(g\) 是从 X 到 \(\mathbf{C}\) 中的连续函数。\(\mu\)-本质像（\(g\) 的）是 \(\mu\) 的支撑在 \(g\) 下的像的闭包。

设 \(Y\) 是 \(\mu\) 的支撑。若 \(z \in \mathbf{C}\) 不附着于 \(g(Y)\)，则存在开邻域 \(U\) （\(z\) 的），使得开集 \(g^{-1}(U)\) 不与 Y 相交，从而是局部 \(\mu\)-零测的；这意味着 \(z\) 不属于 \(\mu\)-本质像（\(g\) 的）。

反之，若 \(z \in \mathbf{C}\) 附着于 \(g(Y)\)，则对 z 的每个开邻域 \(U\) ，集 \(g^{-1}(U)\) 是 X 中与 Y 相交的开集。因此它不是 \(\mu\)-零测的，又由于它是开的，它不是局部 \(\mu\)-零测的（INT, IV, p. 172, § 5, n°2，推论 2）。于是 z 属于 g 的 \(\mu\)-本质像。

引理 3. --- 设 \(g\) 是从 X 到 \(\mathbf{C}\) 中的连续函数，S 是它的 \(\mu\)-本质像。假设 S 非空且 \(0 \notin S\) ，并用 \(\delta\) 表示 0 到 S 的距离。则 \(\delta > 0\)。

置 \(h(x) = 1 / g(x)\) （若 \(g(x) \neq 0\)），否则置 \(h(x) = 0\) 。函数 \(h\) 属于 \(\mathcal{L}^{\infty}(\mathrm{X},\mu)\)。设 \(\widetilde{h}\) 是 \(h\) 在 \(\mathrm{L}^{\infty}(\mathrm{X},\mu)\) 中的类。我们有公式 \(\delta^{-1} = \| \widetilde{h}\|_{\infty}\) 。

设 U 是 0 的开邻域，使得开集 \(Z = g^{-1}(U)\) 局部 \(\mu\)-零测，从而是零测的（INT, IV, p. 172, § 5, no. 2，推论 2）。设 Y 是 \(\mu\) 的支撑；我们有 \(Y \subset X - Z\) 。函数 \(h\) 在 \(X - Z\) 上的限制连续且有界，故 \(h \in \mathcal{L}^{\infty}(X, \mu)\) ，且 \(\widetilde{h}\) 在 \(\mathcal{L}^{\infty}(X, \mu)\) 中的范数等于它在 \(X - Z\) 上的限制的范数。此外，对每个 \(\alpha \in R_{+}\) ，\(x \in X - Z\) 中满足 \(|h(x)| > \alpha\) 者所成之集在 \(X - Z\) 中是开的；因此它局部 \(\mu\)-零测当且仅当它不与 Y 相交（INT, IV，同前及 INT, III, p. 66, § 3, no. 2，定义 1）。于是，我们有

\[
\| \widetilde {h} \| _ {\infty} = \sup _ {x \in \mathrm{Y}} \frac {1}{| g (x) |},
\]

由此得

\[
\frac {1}{\| \widetilde {h} \| _ {\infty}} = \inf _ {x \in \mathrm{Y}} | g (x) | = \inf _ {\lambda \in g (\mathrm{Y})} | \lambda | = \inf _ {\lambda \in \overline {{g (\mathrm{Y})}}} | \lambda |
\]

由前一引理，它等于 δ。

